\documentclass[11pt]{article}

\usepackage{amsmath,amssymb,amstext,amsthm}
\usepackage{geometry}
\usepackage{fancyhdr}
\usepackage[pdftex]{graphicx}
\usepackage{xcolor}

\geometry{
  verbose,
  dvips,
  width=430pt, marginparsep=0pt, marginparwidth=0pt,
  top=90pt, headheight=12pt, headsep=20pt, footskip=30pt, bottom=80pt}

\setlength{\parskip}{\medskipamount}
\setlength{\parindent}{0pt}

\linespread{1.05}

\newtheorem{theorem}{Theorem}
\newtheorem{lemma}[theorem]{Lemma}
\newtheorem{proposition}[theorem]{Proposition}
\newtheorem{corollary}[theorem]{Corollary}
\newtheorem{conjecture}[theorem]{Conjecture}
\newtheorem{obs}{Observation}
\newtheorem{clm}{Claim}
\newtheorem{ques}{Question}
\newtheorem{problem}{Problem}

\newtheorem{ex}{Example}


\newcommand{\spc}{\hspace{0.08in}}
\newcommand{\vs}{\vspace*{.1in}}
\newcommand{\E}{\mathbb{E}}
\newcommand{\Prob}{\mathbb{P}}
\newcommand{\edit}[1]{\emph{\textcolor{blue}{#1}}}


\renewenvironment{proof}{{\noindent \textbf{\textit{Proof.}}}}
{\hfill $\Box$ \vspace*{0.1in}}
\newenvironment{proofc}{{\noindent \textit{Proof of Claim.}}}
{\hfill $\Box$}


\begin{document}

\begin{LARGE}\begin{center}\bf 13.1 Vector Functions and Space Curves\end{center}
\end{LARGE}

\vs\vs



\section{Warm Up}

\textbf{Question 1}: What is \emph{domain}? and What is the domain of $\ln(x-2)$? $\sqrt{x^2-4}$? $\sqrt{x^2+4}$? 

\textbf{Problem 1:} Find $\lim{x \to 0}\frac{1-e^x}{x}$

\section{Motivation}

In two-dimension, The tools of calculus are used on functions and curves. Functions are the language of calculus and it is seemingly natural to try to develop similar functions in 3-dimensions. It turns out that having an output of a function being a vector has many interesting consequences.




\section{Definitions \& Theorems}

\begin{enumerate}
\item  In 3-D a \textbf{parametric function} is of the form $x= f(t), y=g(t), z=h(t)$. Where $t$ is our only independent input (parameter) and for each $t$ in our domain, we get a point. Likewise a \textbf{vector function} $r(t)$ takes in our parameter $t$ and outputs a vector where each component is a function. 

\begin{equation*}
    r(t) = \langle f(t), g(t), h(t) \rangle = f(t)i+g(t)j+h(t)k.
\end{equation*}

Since each vector output is in our domain for $t$, the terminal points of our vector functions outputs, traces out a curve in 3-space which has some orientation. This is known as a \textbf{space curve}.

\item We can imagine now taking the \textbf{limit} of a vector function in very natural way.

\begin{equation*}
    \lim_{t\to a}r(t) = \left\langle \lim_{t\to a}f(t), \lim_{t\to a}g(t), \lim_{t\to a}h(t) \right\rangle.
\end{equation*}


\item We say a vector function is \textbf{continuous} at $a$ if the limit of the vector function as $t$ approaches $a$ is equal to the actual vector function value at $a$. 

\begin{equation*}
    \lim_{t \to a}r(t) = r(a).
\end{equation*}
   

\item The vector function $\langle t,t^2,t^3 \rangle$ gives us the \textbf{twisted cubic}. The vector function $\langle (2+\cos{1.5t})\cos(t),(4+\sin{20t})\sin{t}, \cos{20t} \rangle$ gives the \textbf{trefoil knot}.
   
   
   \end{enumerate}



\section{Examples}

\begin{problem}
    Write the component function for the following vector function, find its domain, and find any vector output.

    \begin{equation*}
        r(t) = \left\langle \frac{1}{t}, \sqrt{t+2}, \ln(t-3) \right\rangle.
    \end{equation*}
\end{problem}

\vspace{8.5cm}

\begin{problem}
    Find the $\lim_{t \to 0}r(t)$ where $r(t) = (1+t^5)i+te^{-t}j+\frac{\sin{t}}{t}k$.
\end{problem}

\vspace{8.5cm}

\begin{problem}
    What does the vector function $r(t) = \cos{t}i+\sin{t}j+tk$ look like?
\end{problem}


\vspace{8.5cm}

\begin{problem}
    Find a vector equation and parametric equations for the line segment which joins the point $P(1,3,-2)$ and $Q(2,-1,3)$.
\end{problem}

\newpage

\section{Scratch Work}

\end{document} 